\chapter{Serial COM1}
\label{cap:serial}

\section{Overview}

The \code{src/serial.rs} module implements a subset of the Linux 8250 driver
for the UART 16550 on COM1. The serial port has two roles:

\begin{itemize}
  \item \textbf{debugging}: all text printed to VGA is also sent to COM1;
  \item \textbf{input}: in addition to the PS/2 keyboard, the shell accepts
        commands typed in the serial terminal, which enables automated tests.
\end{itemize}

There are no interrupts, IRQ-based FIFO, \emph{termios}, \emph{tty} layer, or
flow control. Access is done by polling I/O registers.

\section{Hardware}

The UART 16550 on COM1 uses a base of 8 ports starting at \code{0x3F8}:

\begin{table}[H]
\centering
\caption{UART 16550 registers (COM1).}
\label{tab:uartregs}
\begin{tabular}{@{}llll@{}}
\toprule
\textbf{Offset} & \textbf{Name} & \textbf{Read} & \textbf{Write} \\
\midrule
\code{+0} & RBR / THR & Receive Buffer & Transmit Holding \\
\code{+1} & IER / DLM & Interrupt Enable & Divisor Latch High \\
\code{+2} & IIR / FCR & Interrupt ID & FIFO Control \\
\code{+3} & LCR & Line Control & Line Control \\
\code{+4} & MCR & Modem Control & Modem Control \\
\code{+5} & LSR & Line Status & --- \\
\code{+6} & MSR & Modem Status & --- \\
\code{+7} & SCR & Scratch & Scratch \\
\bottomrule
\end{tabular}
\end{table}

The LSR register has the bits used by the driver:

\begin{table}[H]
\centering
\caption{Line Status Register (LSR) bits.}
\label{tab:lsr}
\begin{tabular}{@{}lll@{}}
\toprule
\textbf{Bit} & \textbf{Mask} & \textbf{Meaning} \\
\midrule
0 & \code{0x01} & Data available for reading \\
5 & \code{0x20} & Transmitter Holding Register empty (can send) \\
\bottomrule
\end{tabular}
\end{table}

\section{Initialization}

The implemented configuration is 38400 baud, 8 data bits, no parity, and 1
stop bit (8N1), with FIFO enabled and IRQs disabled.

\begin{table}[H]
\centering
\caption{COM1 initialization sequence.}
\label{tab:uartinit}
\begin{tabular}{@{}llp{7.2cm}@{}}
\toprule
\textbf{Offset} & \textbf{Value} & \textbf{Effect} \\
\midrule
\code{+1} & \code{0x00} & Disables interrupts. \\
\code{+3} & \code{0x80} & Turns on DLAB to access the divisor. \\
\code{+0} & \code{0x03} & Low divisor = 3. \\
\code{+1} & \code{0x00} & High divisor = 0. \\
\code{+3} & \code{0x03} & 8 bits, no parity, 1 stop; turns off DLAB. \\
\code{+2} & \code{0xC7} & FIFO enabled, clears buffers, 14-byte trigger. \\
\code{+4} & \code{0x0B} & DTR, RTS, and OUT2; IRQs disabled. \\
\bottomrule
\end{tabular}
\end{table}

The divisor 3 produces:

\[
\text{baud} = \frac{115200}{3} = 38400.
\]

After configuration, the driver reads the LSR. If the value is \code{0xFF},
it considers that there is no UART on the bus and sets \code{PRESENT = false}.
This heuristic works because a bus without a device generally returns all bits
as 1.

\begin{nota}[Why 38400 and not 115200?]
The divisor 3 is a conservative value widely supported by emulations and simple
cables. The rate does not need to be high for a text terminal. For 115200, the
divisor would be 1; for 9600, 12. The choice does not affect VGA operation.
\end{nota}

\section{Write}

\begin{lstlisting}[caption={Sending one byte through the UART.},label={lst:uartwrite}]
fn write_byte(byte: u8) {
    unsafe {
        while inb(COM1 + 5) & 0x20 == 0 {}
        outb(COM1, byte);
    }
}
\end{lstlisting}

The loop waits for the THR to empty. If there is no UART and the bus returns
\code{0xFF}, bit \code{0x20} is already on and the loop does not hang; the
byte is simply discarded. If a UART is present but stuck, the loop may wait
indefinitely.

\code{write\_str} converts \code{\textbackslash n} into \code{\textbackslash r
\textbackslash n}, because serial terminals generally expect CRLF:

\begin{lstlisting}[caption={Writing a string with newline conversion.},label={lst:uartstr}]
pub fn write_str(text: &str) {
    for byte in text.bytes() {
        if byte == b'\n' {
            write_byte(b'\r');
        }
        write_byte(byte);
    }
}
\end{lstlisting}

\section{Read}

\begin{lstlisting}[caption={Non-blocking read of one byte.},label={lst:uartread}]
pub fn read_byte() -> Option<u8> {
    if !PRESENT.load(Ordering::Relaxed) {
        return None;
    }
    unsafe {
        if inb(COM1 + 5) & 0x01 == 0 {
            return None;
        }
        Some(inb(COM1))
    }
}
\end{lstlisting}

The function is non-blocking: if there is no byte, it returns \code{None}. The
shell calls it after \code{keyboard::read\_char()}, so the serial port works as
an alternative input source.

\section{API}

\begin{table}[H]
\centering
\caption{API of \code{serial.rs}.}
\label{tab:serialapi}
\begin{tabularx}{\linewidth}{@{}lX@{}}
\toprule
\textbf{Function} & \textbf{Description} \\
\midrule
\code{init()} & Configures COM1 at 38400 8N1 and detects presence. \\
\code{read\_byte()} & Returns a received byte or \code{None}. \\
\code{write\_byte(byte)} & Sends a byte, waiting for the THR. Internal. \\
\code{write\_str(text)} & Sends a string, converting \code{\textbackslash n}
to CRLF. \\
\bottomrule
\end{tabularx}
\end{table}

\section{Usage in QEMU and on physical hardware}

In QEMU, the serial port can be exposed in several ways:

\begin{lstlisting}[style=shell,caption={Examples of exposing the serial port in QEMU.},label={lst:qemuserial}]
# Terminal interativo no stdio (usado por `make run`)
qemu-system-x86_64 -cdrom vga-kernel.iso -boot order=d -m 128M -serial stdio

# Socket Unix para testes automatizados
qemu-system-x86_64 -cdrom vga-kernel.iso -boot order=d -m 128M \
  -display none -serial unix:/tmp/vga-serial.sock,server=on,wait=off
\end{lstlisting}

On a physical PC, COM1 appears at \code{0x3F8} and can be accessed with a
USB-serial cable and a terminal at 38400 8N1. The BIOS may disable the port or
remap it; in that case the driver does not find it.

\section{Limitations}

\begin{itemize}
  \item \textbf{No IRQs}: reading depends on polling.
  \item \textbf{No input FIFO}: bytes may be lost if they arrive faster than
        the shell polls.
  \item \textbf{No flow control}: RTS/CTS are not used.
  \item \textbf{No automatic \emph{echo}}: the shell itself echoes the received
        characters.
  \item \textbf{No support for other ports}: COM1 only.
  \item \textbf{Heuristic detection}: \code{LSR == 0xFF} may not be sufficient
        on all chipsets.
\end{itemize}
